\documentclass[11pt]{article}

\usepackage[T1]{fontenc}
\usepackage{lmodern}
\usepackage{amsmath,amssymb,amsthm,mathtools}
\usepackage[margin=1in]{geometry}
\usepackage{microtype}
\usepackage[hidelinks]{hyperref}

\newtheorem{theorem}{Theorem}
\newtheorem{proposition}[theorem]{Proposition}
\newtheorem{corollary}[theorem]{Corollary}
\theoremstyle{remark}
\newtheorem{remark}[theorem]{Remark}

\newcommand{\R}{\mathbb{R}}
\newcommand{\D}{\mathbb{D}}
\newcommand{\eps}{\varepsilon}

\title{Quantitative Endpoint Classification of a Sufficient Condition in a Weighted Bergman Hilbert-Operator Norm Argument}
\author{Xiaojian Zeng\\
PBC School of Finance, Tsinghua University, 100190\\
\texttt{xiaojianzeng@foxmail.com}}
\date{}

\begin{document}
\maketitle

\begin{abstract}
We audit a specific endpoint step in the proof of a recent norm claim for a
generalized Hilbert operator on weighted Bergman spaces.  Put
$A=\alpha+2\geq2$, $h=p-A$, and
$C_A=B(A-1,1+A/2)$.  For the exact left-hand side $L_{A,b}(h)$ of the
sufficient condition in Corollary~3.5 of Li--Li--Pan--Wulan--Zhou, we prove
an explicit second-order remainder estimate
\[
 L_{A,b}(h)=\frac{\pi^2b}{3}
 +\left\{\frac{\pi^2}{3A}-2b\zeta(3)\left(2-\frac1A\right)\right\}h
 +O_{A,b}(h^2).
\]
The displayed error is made fully quantitative below.  It yields the exact
fixed-$b$ local classification
\[
 \exists\eps>0\ \forall h\in(0,\eps):L_{A,b}(h)\leq C_A
 \quad\Longleftrightarrow\quad
 0\leq b<\frac{3C_A}{\pi^2},
\]
for $0\leq b<1$; equality belongs to the failing side.  We also prove that
the auxiliary upper bound used in the source has right endpoint limit $4b$
and give an explicit neighborhood on which it exceeds $C_A$ when
$b>C_A/4$.  Consequently, the Corollary~3.5 route does not establish the
natural uniform-in-$b$ reading of the source theorem.  Since that corollary
is only sufficient, this result is neither a counterexample to the operator
norm identity nor a disproof of the theorem; the operator-level question
remains unresolved.
\end{abstract}

\section{Context and scope}

Hilbert and generalized Hilbert operators on analytic function spaces have a
substantial literature, including work on generalized operators and exact
norm questions on weighted Bergman spaces
\cite{galanopoulos2014generalized,dai2024norm}.  For $\alpha>-1$, the
standard weighted Bergman space $A_\alpha^p$ consists of analytic functions
on $\D$ with
\[
 \|f\|_{A_\alpha^p}^p
  =(\alpha+1)\int_{\D}|f(z)|^p(1-|z|^2)^\alpha\,dA(z),
\]
where $dA$ is normalized area measure.  The generalized Hilbert operator in
\cite{li2026norm} has the integral representation
\[
 \mathcal H_bf(z)=\int_0^1
 \frac{f(t)(1-t)^b}{(1-tz)^{b+1}}\,dt.
\]
Li et al. study, among other parameter regimes, the identity
\begin{equation}
 \|\mathcal H_b\|_{A_\alpha^p\to A_\alpha^p}
 =B\left(\frac{\alpha+2}{p},
 b+1-\frac{\alpha+2}{p}\right),                         \label{eq:norm}
\end{equation}
and their Theorem~3.4 states a near-endpoint conclusion with a threshold
said to depend only on $\alpha$ \cite[Theorem~3.4]{li2026norm}.

Our result concerns only the sufficient condition (3.38) in
\cite[Corollary~3.5]{li2026norm}.  Set
\[
 A=\alpha+2\geq2,\qquad h=p-A>0,
 \qquad C_A=B\left(A-1,1+\frac A2\right).
\]
The source condition is
\begin{equation}
 L_{A,b}(h)\leq C_A,                                      \label{eq:condition}
\end{equation}
where direct substitution in its exact left-hand side gives
\begin{equation}
 L_{A,b}(h)=\frac{Ab+(b+1)h}{h^2}
 \left\{1-\frac{1}{2hB(A/(A+h),2h)}\right\}.             \label{eq:L}
\end{equation}
The source domain also imposes $0\leq b<A/(A+h)$.  Thus a fixed $b$ has
admissible points arbitrarily close to $h=0$ exactly when $0\leq b<1$.
Corollary~3.5 additionally requires $p<p_\alpha$, where
\[
 p_\alpha=\frac{3\alpha}{4}+2+
 \frac14\sqrt{9\alpha^2+40\alpha+48}>A+\frac12.
\]
All neighborhoods below are capped at $h\leq1/4$, so this latter restriction
is automatic.

The closest prior statement is precisely Corollary~3.5: it proves
\eqref{eq:norm} when \eqref{eq:condition} holds.  The delta here is an
explicit endpoint expansion, a sharp fixed-$b$ local classification
including the equality case, and the resulting diagnosis of the particular
proof route.  No necessity claim for \eqref{eq:condition} is made.

\section{The auxiliary endpoint bound}

The proof of Theorem~3.4 in \cite{li2026norm} uses the function
\begin{equation}
 G_{\alpha,b}(p)=
 \frac{p(b+1)-\alpha-2}{p-\alpha-2}
 \frac{2\left[1-\left((\alpha+2)/p\right)^2\right]}
 {\, (\alpha+2)/p+2(p-\alpha-2)\,}.                     \label{eq:Gsource}
\end{equation}

\begin{proposition}\label{prop:G}
Let $A\geq2$ and $0\leq b<1$.  With $p=A+h$,
\[
 \lim_{h\downarrow0}G_{A-2,b}(A+h)=4b.
\]
Moreover, let $C_A=B(A-1,1+A/2)$ and suppose $b>C_A/4$.  Define
\[
 M(A,b)=4A^2b+9Ab+2A+5b+1,
 \qquad K(A,b)=\frac{2M(A,b)}{A^2},
\]
and $\Delta=4b-C_A$.  Then
\begin{equation}
 \eps_G=\min\left\{\frac12,\frac{A(1-b)}{2b},
                  \frac{\Delta}{2K(A,b)}\right\}>0      \label{eq:epsG}
\end{equation}
has the following property: every $0<h<\eps_G$ is admissible,
$b<A/(A+h)$, and satisfies
$G_{A-2,b}(A+h)>C_A$.
\end{proposition}

\begin{proof}
Cancellation in \eqref{eq:Gsource} gives
\begin{equation}
 G_{A-2,b}(A+h)=
 \frac{2(Ab+(b+1)h)(2A+h)}
 {(A+h)(A+2Ah+2h^2)}.                                   \label{eq:Grational}
\end{equation}
The denominator tends to $A^2$ and the numerator to $4A^2b$, proving the
limit.  More quantitatively,
\begin{equation}
 G_{A-2,b}(A+h)-4b
 =-\frac{2hQ(A,b,h)}{(A+h)(A+2Ah+2h^2)},                 \label{eq:Gdiff}
\end{equation}
where
\[
 Q(A,b,h)=4A^2b+8Abh-Ab-2A+4bh^2-bh-h.
\]
For $0<h\leq1$, the denominator in \eqref{eq:Gdiff} is at least $A^2$
and $|Q(A,b,h)|\leq M(A,b)$.  Hence
\begin{equation}
 |G_{A-2,b}(A+h)-4b|\leq K(A,b)h.                       \label{eq:Glinear}
\end{equation}
The second entry in \eqref{eq:epsG} implies
$bh<A(1-b)$, hence $b<A/(A+h)$.  The third entry and
\eqref{eq:Glinear} give
\[
 G_{A-2,b}(A+h)>4b-\frac\Delta2
  =C_A+\frac\Delta2>C_A.
\]
\end{proof}

The conclusion $G>C_A$ says only that this particular upper estimate no
longer closes; an upper bound larger than $C_A$ cannot by itself show that
the exact quantity $L_{A,b}$ exceeds $C_A$.  The next section determines
the latter question directly.

\section{Quantitative expansion of the exact condition}

We use the classical power series for $\log\Gamma(1+z)$; a standard
reference is DLMF~\S5.7(i) \cite{dlmf}.

\begin{theorem}[Explicit endpoint remainder]\label{thm:remainder}
Let $A\geq2$, $0\leq b<1$, and $0<h\leq1/4$.  Define
\begin{align*}
 c_2&=\frac{\pi^2}{3A},\\
 c_3&=-\frac{\pi^2}{3A^2}
       +\zeta(3)\left(\frac2{A^2}-\frac4A\right),\\
 K_\Phi&=\frac{\pi^2}{3A^3}+\frac{\zeta(3)}2+\frac{897}{28},\\
 P&=c_2+\frac{|c_3|}{4}+\frac{K_\Phi}{16},\\
 E&=\frac12 e^{P/16}P^2,\qquad K_N=K_\Phi+E,
\end{align*}
and
\begin{align}
 D(A,b)&=\frac{\pi^2}{3A}
       -2b\zeta(3)\left(2-\frac1A\right),                \label{eq:D}\\
 K_L(A,b)&=AbK_N+(b+1)|c_3|+\frac{b+1}{4}K_N.             \label{eq:KL}
\end{align}
Then
\begin{equation}
 \left|L_{A,b}(h)-\frac{\pi^2b}{3}-D(A,b)h\right|
 \leq K_L(A,b)h^2.                                      \label{eq:mainbound}
\end{equation}
\end{theorem}

\begin{proof}
Put
\[
 d=-\frac{h}{A+h},\qquad u=2h,\qquad F(z)=\log\Gamma(1+z).
\]
The beta--gamma identity gives
\begin{equation}
 2hB\left(\frac{A}{A+h},2h\right)=e^{\Phi(h)},
 \qquad
 \Phi(h)=F(d)+F(u)-F(d+u).                               \label{eq:Phi}
\end{equation}
and therefore
\begin{equation}
 L_{A,b}(h)=\frac{Ab+(b+1)h}{h^2}\bigl(1-e^{-\Phi(h)}\bigr).
                                                                    \label{eq:Lphi}
\end{equation}
For $0<h\leq1/4$, we have $|d|\leq h/2$, $u\leq1/2$, and
$0<d+u\leq2h$.  Thus the absolutely convergent series
\begin{equation}
 F(z)=-\gamma z+\sum_{n=2}^{\infty}
       \frac{(-1)^n\zeta(n)}{n}z^n,\qquad |z|<1,         \label{eq:loggamma}
\end{equation}
may be applied to all three terms in \eqref{eq:Phi}.

The linear terms cancel.  The quadratic term is
\[
 -\zeta(2)du=\frac{\pi^2h^2}{3(A+h)}
 =\frac{\pi^2}{3A}h^2-\frac{\pi^2}{3A^2}h^3+R_2,
 \qquad |R_2|\leq\frac{\pi^2}{3A^3}h^4.
\]
The cubic term is
\[
 \zeta(3)(d^2u+du^2)
 =\zeta(3)h^3\left(\frac2{(A+h)^2}-\frac4{A+h}\right).
\]
For $s\geq2$, the derivative of $2s^{-2}-4s^{-1}$ has absolute value at
most $1/2$.  Replacing $A+h$ by $A$ in the last display therefore creates
an error at most $\zeta(3)h^4/2$.

For the remaining terms, $\zeta(n)<2$ and $2/n\leq1/2$ for $n\geq4$,
so
\begin{align*}
 |R_{\geq4}|
 &\leq\frac12\sum_{n=4}^{\infty}
   \left\{(h/2)^n+2(2h)^n\right\}\\
 &\leq\left(\frac1{28}+32\right)h^4
   =\frac{897}{28}h^4.
\end{align*}
Combining these estimates yields
\begin{equation}
 |\Phi(h)-c_2h^2-c_3h^3|\leq K_\Phi h^4.                 \label{eq:Phibound}
\end{equation}

In particular, $|\Phi(h)|\leq Ph^2\leq P/16$.  The elementary inequality
\[
 |1-e^{-y}-y|\leq\frac12e^{|y|}y^2\qquad(y\in\R)
\]
and \eqref{eq:Phibound} now imply
\begin{equation}
 1-e^{-\Phi(h)}=c_2h^2+c_3h^3+R_N(h),
 \qquad |R_N(h)|\leq K_Nh^4.                            \label{eq:Nbound}
\end{equation}
Multiplying \eqref{eq:Nbound} by the prefactor in \eqref{eq:Lphi}, the
constant and linear coefficients are
\[
 Abc_2=\frac{\pi^2b}{3},
 \qquad
 Abc_3+(b+1)c_2
 =\frac{\pi^2}{3A}-2b\zeta(3)\left(2-\frac1A\right).
\]
The remaining terms have absolute value at most
\[
 \left\{AbK_N+(b+1)|c_3|+\frac{b+1}{4}K_N\right\}h^2,
\]
which is exactly \eqref{eq:mainbound}.
\end{proof}

\section{Sharp local classification}

\begin{theorem}[Endpoint classification]\label{thm:classification}
Let $A\geq2$, $0\leq b<1$, and $C_A=B(A-1,1+A/2)$.  There is an
admissible punctured right neighborhood on which the source condition
\eqref{eq:condition} holds if and only if
\begin{equation}
 0\leq b< b_*(A),\qquad b_*(A)=\frac{3C_A}{\pi^2}.        \label{eq:bstar}
\end{equation}
At $b=b_*(A)$ the condition fails throughout an explicit punctured right
neighborhood.
\end{theorem}

\begin{proof}
The beta integral gives
\[
 0<C_A=\int_0^1t^{A-2}(1-t)^{A/2}\,dt\leq\frac12,
\]
and $\pi^2>9$, so $0<b_*(A)<1$.  To keep every point admissible, define
\begin{equation}
 \rho(A,b)=
 \begin{cases}
  1/4,&b=0,\\[2mm]
  \min\{1/4,A(1-b)/(2b)\},&b>0.
 \end{cases}                                             \label{eq:rho}
\end{equation}
For $0<h<\rho(A,b)$, one has $b<A/(A+h)$ and $p=A+h<p_\alpha$.

First suppose $b\neq b_*(A)$ and set
\[
 J(A,b)=|D(A,b)|+\frac14K_L(A,b),
\]
where $D$ and $K_L$ are given in \eqref{eq:D}--\eqref{eq:KL}.  The explicit
choice
\begin{equation}
 \eps_{\neq}=\min\left\{\rho(A,b),
 \frac{|C_A-\pi^2b/3|}{2J(A,b)}\right\}                 \label{eq:epsneq}
\end{equation}
is positive.  Theorem~\ref{thm:remainder} implies, for
$0<h<\eps_{\neq}$,
\[
 \left|L_{A,b}(h)-\frac{\pi^2b}{3}\right|
 \leq J(A,b)h<\frac12\left|C_A-\frac{\pi^2b}{3}\right|.
\]
Consequently $L_{A,b}(h)<C_A$ for $b<b_*(A)$, while
$L_{A,b}(h)>C_A$ for $b>b_*(A)$.

It remains to decide equality, where the constant term is exactly $C_A$.
Substitution of \eqref{eq:bstar} into \eqref{eq:D} shows that
$D(A,b_*)>0$ is equivalent to
\begin{equation}
 \pi^4>18C_A\zeta(3)(2A-1).                              \label{eq:Dequiv}
\end{equation}
Now $C_A\leq1/(A-1)$, the integral test gives $\zeta(3)<5/4$, and
$(2A-1)/(A-1)\leq3$ for $A\geq2$.  Hence
\[
 18C_A\zeta(3)(2A-1)
 \leq\frac{45}{2}\frac{2A-1}{A-1}
 \leq\frac{135}{2}<81<\pi^4,
\]
which proves \eqref{eq:Dequiv}.  Therefore
\begin{equation}
 \eps_{=}=\min\left\{\rho(A,b_*),
 \frac{D(A,b_*)}{2K_L(A,b_*)}\right\}>0                \label{eq:epseq}
\end{equation}
satisfies, for $0<h<\eps_{=}$,
\[
 L_{A,b_*}(h)-C_A
 \geq D(A,b_*)h-K_L(A,b_*)h^2
 >\frac12D(A,b_*)h>0.
\]
This assigns equality to the failing side and completes the classification.
\end{proof}

\section{Implication for the source proof and limitations}

Equation~(3.41) of \cite{li2026norm} is the quantity
$G_{\alpha,b}(p)$ in \eqref{eq:Gsource}.  Although its second factor tends
to zero at $p=\alpha+2$, its first factor is singular there; the product
tends to $4b$, not to zero.  Proposition~\ref{prop:G} makes the resulting
failure of that upper-bound argument quantitative.  Theorem
\ref{thm:classification} goes further: for every fixed
$b\in[b_*(A),1)$, the exact sufficient condition (3.38) itself is false at
all sufficiently close admissible points to the right of the endpoint.

Theorem~3.4 of \cite{li2026norm} places $p$ in a hypothesis before later
quantifying over $p$.  Under its natural reading---a threshold depending
only on $\alpha$, uniform over admissible $b$---the proof explicitly ends
by invoking Corollary~3.5.  Since the corollary's hypothesis fails locally
for the nonempty range $b\in[b_*(A),1)$, that route does not prove this
uniform reading.  The result above does not determine whether a different
argument proves the norm identity, whether a valid threshold must depend on
$b$, or whether another additional hypothesis suffices.

Most importantly, condition \eqref{eq:condition} is sufficient, not known
to be necessary.  Its failure supplies no function witnessing a violation
of \eqref{eq:norm}.  For comparison, an actual finite-dimensional
matrix-norm disproof by Zeng--Liu--Ratnavelu exhibits explicit matrices
and separates the attained and conjectured constants by exact rational
inequalities \cite{zeng2026counterexample}.  No analogous operator-level
witness or necessary-condition argument is supplied here.  Accordingly,
the status of the source's operator-norm assertion is left unresolved.

\section*{Computational and AI assistance}

Two accompanying Python scripts were used for error detection.  SymPy
1.14 \cite{meurer2017sympy} checked the rational simplification, endpoint
limits, and symbolic coefficients.  A separate mpmath 1.3.0 computation at
90 decimal digits compared the beta and log-gamma representations and
tested the remainder bound at 144 finite parameter points; the largest
observed error-to-bound ratio was approximately $0.187283$.  These finite
checks are supplementary and prove none of the quantified statements above.
OpenAI Codex assisted with proof exploration, source comparison, adversarial
checking, and manuscript preparation.  The author checked the mathematical
arguments and is responsible for the final text and for compliance with any
venue's authorship and disclosure policies.

\bibliographystyle{plain}
\bibliography{references}

\end{document}
